% ============================================================================
% instantiations/lu-university-of-luxembourg/set.tex — INSTITUTION SET (Layer 2)
% ----------------------------------------------------------------------------
% LUXEMBOURG REFERENCE INSTANTIATION — University of Luxembourg.
%
% This is the canonical, named concrete edition of the SUIT suite. It binds the
% generic role-phrases to the real Luxembourg / University of Luxembourg surface
% forms. Every binding below is the University-of-Luxembourg surface form for its
% localization variable (see shared/localization-defaults.tex for the variable
% set and Layer-0 generic defaults). A concrete instantiation legitimately names
% Luxembourg and the University of Luxembourg: that is its purpose. It is NOT the
% generic body.
%
% Localization model:
%     localization.sty                         (\setloc / \loc machinery)
%       Layer 0: shared/localization-defaults.tex   (generic role phrases)
%         Layer 1: categories/EU-PUB-RI.tex          (European public RI class)
%           Layer 2: THIS FILE                        (concrete LU/UL values)
%             document body
%
% NOTE ON GLOSSARY: the SUIT glossary (shared/glossary-defs.tex) is
% institution-neutral — it defines NO UL-named acronym entries (no \gls{govcertlu},
% \gls{anssilu}, \gls{monarc}, \gls{psilu}, ...). The matrix surface forms carry
% such \gls{} wrappers because they came from the original UL-suite glossary; we
% therefore bind PLAIN-TEXT values here (calling an undefined \gls key would fail
% the build). eduGAIN is the only LU-relevant glossary key that exists upstream.
% ============================================================================

% --- 1. CHOSEN CATEGORY ARCHETYPE (Layer 1) ---------------------------------
% University of Luxembourg: a public, research-intensive university established
% under the law of an EU/EEA Member State — the worked EU-PUB-RI archetype lists
% it as a representative member.
\newcommand{\instcategory}{EU-PUB-RI}
\input{\catdir/\instcategory}% -> categories/EU-PUB-RI.tex

% ============================================================================
% 2. COUNTRY-SPECIFIC NAME KEYS (left at default by the EU-PUB-RI archetype)
% ----------------------------------------------------------------------------
% The archetype binds the structural and supranational-regulatory keys and
% leaves these per-country NAME keys for the Member-State edition to fill.
% LU bindings from the matrix (ANCHOR-NREN, ANCHOR-DPA, ANCHOR-GOV-CSIRT).
% ============================================================================
\setloc{nren-name}{RESTENA (the Luxembourg national research and education network)}
\setloc{national-dpa}{the CNPD (Commission Nationale pour la Protection des Donn\'ees, the Luxembourg national data-protection authority)}
\setloc{national-gov-csirt}{GOVCERT.LU (the Luxembourg governmental CSIRT)}

% ============================================================================
% 3. INSTITUTION IDENTITY (named LU edition)
% ----------------------------------------------------------------------------
% ANCHOR-INSTITUTION-NAME, ANCHOR-INSTITUTION-TYPE-PROFILE, ANCHOR-INSTITUTION-
% MACRO-SET, and the governance ladder (ANCHOR-EXECUTIVE-LEADERSHIP,
% ANCHOR-GOVERNING-BOARD, ANCHOR-ACADEMIC-BODIES).
% ============================================================================
\setloc{institution-name}{the University of Luxembourg}
\setloc{institution-macro-set}{the University of Luxembourg}
\setloc{institution-type}{a public, research-intensive (R1) university established under the law of the Grand Duchy of Luxembourg}

% Governance ladder (the named University Bodies of the University of Luxembourg):
\setloc{executive-leadership}{the Rectorate (the Rector and the central executive team)}
\setloc{governing-board}{the Board of Governors (the Conseil de Gouvernance)}
\setloc{academic-governing-bodies}{the academic deliberative bodies (the Faculty Councils and the University Council)}
\setloc{staff-representation-body}{the staff delegation (d\'el\'egation du personnel)}

% Approving academic units for the internal-approval byline (the original UL
% surface form: the named STM faculty and Computer Science department domains):
\setloc{approving-academic-units}{the Faculty of Science, Technology and Medicine and the Department of Computer Science of the University of Luxembourg}

% ============================================================================
% 4. CONCRETE LU/UL VALUES (institution units, pilots, people, network tier)
% ----------------------------------------------------------------------------
% These re-bind keys the archetype set to a generic EU phrase, pinning them to
% the original UL surface forms recorded in the matrix.
% ============================================================================

% --- Institutional units & pilots -------------------------------------------
% ANCHOR-CS-UNIT, ANCHOR-SECURITY-RESEARCH-CENTRE, ANCHOR-INSTITUTION-HPC,
% ANCHOR-PILOT-ENTITIES.
\setloc{cs-research-unit}{the Department of Computer Science (DCS)}
\setloc{security-research-centre}{the SnT (the Interdisciplinary Centre for Security, Reliability and Trust)}
\setloc{institution-hpc-facility}{the University's HPC facility and team (ULHPC)}
\setloc{pilot-entities}{the University's HPC team (ULHPC) together with the Department of Computer Science (DCS) and the SnT (the Interdisciplinary Centre for Security, Reliability and Trust)}

% --- Named author & affiliation ---------------------------------------------
% ANCHOR-AUTHOR-NAME, ANCHOR-AUTHOR-AFFILIATION.
\setloc{author-name}{Nicolas Guelfi}
% Rendered as the OPTIONAL trailing byline clause via \loc{author-affiliation}
% (see localization-defaults.tex): a leading space + a closing full stop yield a
% clean second sentence after the working-group attribution.
\setloc{author-affiliation}{ Department of Computer Science, University of Luxembourg, Luxembourg (nicolas.guelfi@uni.lu). The suite was originally developed for the University of Luxembourg.}

% --- Title-block keys (consumed by \suitmetablock on the four title pages) ---
% Edition facts for the LU edition; bound explicitly so the edition is immune
% to Layer-0 default changes and never renders the Layer-0 approval-date
% placeholder (audit 260609-2153, INST-47f3e9e5).
\setloc{doc-version}{2.1}
\setloc{doc-visibility}{INTERNAL}
\setloc{doc-approval-authority}{\loc{approving-academic-units}}
\setloc{doc-approval-date}{May 7, 2026}

% --- Staffing / FTE profile -------------------------------------------------
% ANCHOR-FTE-ENVELOPE: the original UL framing (R1 institution, small operating
% team, governance bodies setting excellence as the strategic objective).
\setloc{staffing-profile}{the University of Luxembourg, whose governance bodies set excellence as a strategic objective and which sustains the identity needs of an R1 institution within a small operating team (indicatively two to four engineers shared across other duties)}

% --- Research & education network tier (RESTENA + GEANT) --------------------
% ANCHOR-NREN-CSIRT, ANCHOR-IDENTITY-FEDERATION, ANCHOR-RESEARCH-BACKBONE-SCALE.
\setloc{nren-csirt}{RESTENA-CSIRT (the CSIRT of the Luxembourg NREN RESTENA)}
\setloc{identity-federation-profile}{the RESTENA-operated Luxembourg academic identity federation, interconnected through the eduGAIN interfederation (with collaboration authorisation patterns developed in G\'EANT's eduTEAMS service)}
\setloc{research-backbone-scale}{the G\'EANT pan-European research backbone (40+ NRENs serving some 50 million users in 10\,000+ institutions, with multiple-wavelength backbone capacity at 400\,Gbps), reached from Luxembourg through RESTENA with campus connections in the 10--100\,Gbps range}

% --- High-performance computing (national/supranational tier) ----------------
% ANCHOR-EUROHPC-SYSTEM: the national/regional supercomputer is MeluXina.
\setloc{national-hpc-facility}{MeluXina (the Luxembourg EuroHPC JU supercomputer)}

% --- Data custody / sovereignty optimum -------------------------------------
% ANCHOR-ONPREM-KEY-CUSTODY: on-premises storage, institutional key custody,
% no third-country disclosure regime (the structural sovereignty optimum).
\setloc{data-custody-sovereignty-profile}{on-premises storage with the trust anchor held in an HSM under University custody and no third-country disclosure exposure (the structural data-sovereignty optimum), keeping the core estate within the EU/EEA}

% --- Incumbent proprietary identity suite -----------------------------------
% ANCHOR-PROPRIETARY-IDP: the incumbent commercial identity suite is Microsoft
% Entra ID with Conditional Access.
\setloc{incumbent-identity-suite-vendor}{Microsoft Entra ID with Conditional Access}

% --- Internal operational contact -------------------------------------------
% ANCHOR-INTERNAL-EMERGENCY-CONTACT: the UL internal emergency number is 9911.
\setloc{internal-emergency-incident-contact}{the University's internal emergency / incident-reporting contact (emergency number 9911), reporting to the ISO/CISO}

% ============================================================================
% 5. NATIONAL REGULATORY STACK, AUTHORITIES AND INSTRUMENTS (Luxembourg)
% ----------------------------------------------------------------------------
% The EU-PUB-RI archetype binds the supranational instruments (NIS2, GDPR,
% eIDAS 2, AI Act). Here we pin the Luxembourg national transpositions,
% authorities, strategy, methods and the named legal instruments — every value
% is the original UL surface form from the matrix.
% ============================================================================

% --- Composite regime & supranational anchors (named for the LU edition) -----
% ANCHOR-REGULATORY-STACK, ANCHOR-NIS-REGIME, ANCHOR-DATA-PROTECTION-REGIME,
% ANCHOR-EIDAS-REGIME, ANCHOR-AI-REGIME, ANCHOR-TRANSFER-REGIME,
% ANCHOR-SANCTION-CEILING.
\setloc{regulatory-regime}{the European Union regime as it binds Luxembourg, instantiated through the national transpositions of the Grand Duchy of Luxembourg}
\setloc{cybersecurity-baseline-regime}{the NIS2 Directive (EU) 2022/2555 as transposed in Luxembourg}
\setloc{data-protection-regime}{the EU General Data Protection Regulation (Regulation (EU) 2016/679, GDPR), supervised in Luxembourg by the CNPD}
\setloc{eid-trust-services-regime}{the EU eIDAS\,2 trust-services regime (Regulation (EU) No~910/2014 as amended by Regulation (EU) 2024/1183)}
\setloc{ai-regulation}{the EU AI Act (Regulation (EU) 2024/1689)}
\setloc{cross-border-transfer-regime}{the GDPR Chapter~V cross-border data-transfer regime (post-Schrems~II)}
\setloc{nis2-sanction-ceiling}{up to EUR~10 million or 2\% of annual turnover}

% --- National transposition instruments -------------------------------------
% ANCHOR-NIS2-TRANSPOSITION (PL 8364), ANCHOR-CER-TRANSPOSITION (PL 8307),
% ANCHOR-EQUAL-TREATMENT-LAW (Law of 28 Nov 2006), ANCHOR-STATE-INFOSEC-POLICY
% (PSI-LU), ANCHOR-NATIONAL-CYBER-STRATEGY (NCSS III), ANCHOR-RISK-METHOD
% (MONARC), ANCHOR-HCPN-FOUNDING-LAW (loi modifiee du 23 juillet 2016).
\setloc{nis2-national-transposition}{the Luxembourg transposition of the NIS2 Directive (projet de loi 8364, deposited at the Chambre des D\'eput\'es on 13 March 2024)}
\setloc{cer-national-transposition}{the Luxembourg transposition of the CER Directive (EU) 2022/2557 (projet de loi 8307, deposited at the Chambre des D\'eput\'es on 1 September 2023)}
\setloc{national-equal-treatment-law}{the Luxembourg Law of 28 November 2006 (transposing Directive 2000/78/EC into national labour law, integrating the obligation of reasonable adjustments)}
\setloc{national-public-sector-isp}{PSI-LU (the Politique g\'en\'erale de la s\'ecurit\'e de l'information de l'\'Etat luxembourgeois, approved 13 July 2018)}
\setloc{national-cyber-strategy}{NCSS~III (the third National Cybersecurity Strategy of Luxembourg, 2018, Implementation Objective~7)}
\setloc{national-risk-method}{MONARC (the Luxembourg-developed ISO/IEC~27005-aligned risk-analysis method, developed by NC3 under the Luxembourg House of Cybersecurity)}
% Reference-only binding (no \loc{} consumption site; see the note on the
% reference-only authority bindings below): the HCPN founding statute.
\setloc{national-cyber-commissariat-founding-law}{the loi modifi\'ee du 23 juillet 2016 establishing the Haut-Commissariat \`a la Protection nationale (HCPN), parent of ANSSI Luxembourg and host of GOVCERT.LU}

% --- National authorities & bodies ------------------------------------------
% ANCHOR-NATIONAL-INFOSEC-AGENCY (ANSSI-LU), ANCHOR-NIS2-SUPERVISORY-AUTH (ILR),
% ANCHOR-NATIONAL-SPOC-HCPN (HCPN), ANCHOR-CYBER-COMPETENCE-CENTRE (NC3 /
% Luxembourg House of Cybersecurity), ANCHOR-NATIONAL-NOTIFICATION-PLATFORM
% (SERIMA), ANCHOR-NATIONAL-LEGISLATIVE-CHAMBER (Chambre des Deputes),
% ANCHOR-NATIONAL-EXECUTIVE-DEPARTMENT (Ministere d'Etat).
\setloc{national-infosec-agency}{ANSSI Luxembourg (ANSSI-LU, the national Agency for Information Systems Security, a component of HCPN and distinct from the French ANSSI)}
\setloc{nis2-supervisory-authority}{the ILR (the Institut Luxembourgeois de R\'egulation, the Luxembourg NIS2 supervisory authority)}
\setloc{national-spoc-authority}{the HCPN (the Haut-Commissariat \`a la Protection nationale, the Luxembourg NIS2 national single point of contact)}
\setloc{national-cyber-competence-centre}{NC3 (the National Cybersecurity Competence Center) under the Luxembourg House of Cybersecurity}
\setloc{national-notification-platform}{SERIMA (the Luxembourg SEcurity RIsk MAnagement incident-notification platform, operational since 2 May 2025)}

% Reference-only bindings: the following three name the LU legislative chamber,
% the executive department and the HCPN founding statute. The generic suite has
% no \loc{} consumption site for them (they are administrative back-references,
% not surfaces of the rendered argument), and the HCPN/ANSSI-LU substance they
% support is already carried through \loc{national-spoc-authority},
% \loc{national-infosec-agency} and \loc{national-gov-csirt}. They are retained
% as authoritative parameterize anchors for any edition or tooling that does
% cite them; they intentionally do not render in the four SUIT documents.
\setloc{national-legislature}{the Chambre des D\'eput\'es du Grand-Duch\'e de Luxembourg}
\setloc{national-executive-department}{the Minist\`ere d'\'Etat (under the authority of the Prime Minister)}

% --- Essential-entity classification basis ----------------------------------
% ANCHOR-ESSENTIAL-ENTITY-CLASSIFICATION: UL classified by the ILR as an
% essential entity on the basis of its status as a public institution
% (public-administration-entity route, regardless of size).
\setloc{essential-entity-classification-basis}{the University of Luxembourg's classification by the ILR as an essential entity under the NIS2 essential-entity regime, on the basis of its status as a public institution (the public-administration-entity route, which applies regardless of size)}

% ============================================================================
% 6. NOTE ON THE REMAINING UNCONSUMED BINDINGS (substance-in-prose orphans)
% ----------------------------------------------------------------------------
% Some keys bound above have no \loc{} consumption site in the four SUIT
% documents AS RENDERED, yet are deliberately retained (recounted by grep on
% 2026-06-10, audit 260609-2153, INST-4fbe80b2):
%
%   - Declared orphans (no consumption site today): regulatory-regime,
%     eid-trust-services-regime, data-custody-sovereignty-profile,
%     institution-hpc-facility, academic-governing-bodies,
%     internal-emergency-incident-contact, institution-macro-set, author-name,
%     essential-entity-classification-basis, and the reference-only authority
%     bindings noted above (national-legislature, national-executive-department,
%     national-cyber-commissariat-founding-law).
%
% These are NOT regressions: the LU-applicable substance they carry is already
% present in the generic prose (on-prem HSM key custody, the EU/EEA sovereignty
% posture, the internal emergency contact pattern, etc.), so the round-trip
% already closes for them through the body text. They remain bound here as
% authoritative parameterize anchors so that (a) other editions/archetypes that
% DO add a consumption site inherit the correct LU value, and (b) the matrix's
% parameterize intent is recorded against a real value rather than dropped.
% All other bound keys have at least one \loc{} consumption site in the
% generic suite and render in the LU edition --- including the retrofit set
% (national-dpa, national-gov-csirt, national-public-sector-isp,
% cer-national-transposition, national-equal-treatment-law, author-affiliation,
% approving-academic-units, national-spoc-authority,
% national-cyber-competence-centre, national-risk-method, nren-csirt,
% nis2-sanction-ceiling, executive-leadership), the six formerly listed here as
% orphans and consumed since (institution-type, data-protection-regime,
% cybersecurity-baseline-regime, identity-federation-profile,
% incumbent-identity-suite-vendor, research-backbone-scale), and the newer
% staff-representation-body and doc-* title-block keys.
% ============================================================================

\endinput
